\documentclass[11pt]{article}

\usepackage[utf8]{inputenc}
\usepackage[T1]{fontenc}
\usepackage{microtype}
\usepackage[margin=1in]{geometry}
\usepackage{amsmath,amssymb,amsthm}
\usepackage{mathtools}
\usepackage{bm}
\usepackage{booktabs}
\usepackage{graphicx}
\usepackage{xcolor}
\usepackage[round]{natbib}
\usepackage[colorlinks=true,allcolors=blue]{hyperref}
\usepackage{enumitem}

\newtheorem{definition}{Definition}
\newtheorem{assumption}{Assumption}

% Shorthands in the style of the companion papers.
\newcommand{\R}{R}                       % representation-valued random variable
\newcommand{\T}{T}                       % target / linguistic property
\newcommand{\rade}{\mathfrak{R}}         % Rademacher complexity
\newcommand{\hypclass}{\mathcal{H}}      % hypothesis class
\newcommand{\complexity}{c}              % complexity knob
\newcommand{\fit}{\mathrm{fit}}
\newcommand{\IIA}{\mathrm{IIA}}
\newcommand{\Exp}{\mathbb{E}}
\newcommand{\bigO}{\mathcal{O}}

\title{\textbf{The Complexity--Spuriousness Trade-off in Probing:\\
Toward a Unified Capacity Measure for Diagnostic\\ and Causal Interpretability}}

\author{Tiago Pimentel\\ \small ETH Z\"urich \\ \small \texttt{tiago.pimentel@inf.ethz.ch}}

\date{\today}

\begin{document}
\maketitle

\begin{abstract}
\noindent
Two recent results paint an uncomfortable picture of how we interpret neural
networks. On the diagnostic side, \citet{pimentel2020information} show that, under
an information-theoretic operationalisation, a sufficiently powerful probe will
\emph{always} recover a target property---making the success of probing
uninformative about the representation itself. On the causal side,
\citet{sutter2025nonlinear} show that, with a sufficiently complex alignment map,
\emph{any} algorithm can be cast as a distributed abstraction of \emph{any}
network---making causal abstraction vacuous without constraints. Both results are
\emph{population-level} statements about a limit of infinite probe capacity. In
practice, however, we never operate in that limit: we fit a probe or an alignment
map on a finite sample, and the danger is the \emph{finite-sample optimism} that
this fitting introduces---a propensity to find structure that is not there. In
this technical report we argue that this propensity is, for both techniques, an
instance of the same object: the \textbf{Rademacher complexity} of the
interpretability method's hypothesis class, evaluated on the representations. We
make this unification precise, separate it cleanly from the population-level
vacuity that the two theorems describe, and propose a concrete experimental
protocol---building directly on an existing implementation of the hierarchical
equality task---to measure how this \emph{spurious-fitting capacity} scales with
probe complexity for diagnostic probing and for distributed alignment search
(DAS). Our central question is not \emph{whether} both methods fail in the limit
(they do), but \emph{how gracefully} each approaches that limit, and which
structural constraints govern the rate.
\end{abstract}

\section{Motivation}\label{sec:motivation}

A probe is a supervised model trained to predict some property $\T$ from a
network's representations $\R$; if the probe succeeds, we are tempted to conclude
that the representations \emph{encode} the property
\citep{alain2017understanding,hewitt2019structural,belinkov2022probing}.
\citet{pimentel2020information} recast this practice information-theoretically:
probing estimates the mutual information $\mathrm{I}(\T;\R)$, and the
best estimate is obtained by the \emph{best-performing} probe one can train,
regardless of its complexity. The corollary is uncomfortable: under their
\autoref{sec:setup} assumptions, a contextual representation contains exactly as
much information about a word-level property as the underlying sentence does, so a
powerful enough probe recovers the property for reasons that need have nothing to
do with how---or whether---the network uses it.

Causal probing was meant to escape this critique by intervening rather than
merely decoding. Distributed alignment search \citep[DAS;][]{geiger2024finding}
learns a (possibly non-linear) alignment map $\phi$ onto the hidden space, performs
interchange interventions in the aligned basis, and measures the
\emph{interchange-intervention accuracy} (IIA) with which the network reproduces a
target algorithm's counterfactual behaviour.
\citet{sutter2025nonlinear} show that this escape is illusory: once the linearity
constraint on $\phi$ is lifted, \emph{any} algorithm can be aligned to \emph{any}
network with near-perfect IIA---even randomly initialised networks that cannot
solve the task. They call this the \textbf{non-linear representation dilemma}:
lifting the linearity constraint removes any principled way to trade off the
alignment map's complexity against its accuracy.

Read together, the two results describe \emph{the same failure with two faces}.
In both, an interpretability method reports a score---probe accuracy, or
IIA---that is optimised over a family of extraction functions, and in both, that
score becomes uninformative once the family is expressive enough. The practitioner
is left without a principled way to choose the family's complexity. This report
asks whether the two methods are not just analogous but \emph{measurable on a
common axis}: can we define a single notion of capacity, in the spirit of
Rademacher complexity \citep{bartlett2002rademacher} or VC dimension
\citep{vapnik1998statistical}, that quantifies each method's propensity to find
spurious structure, and study how that propensity grows?

\section{Setup and notation}\label{sec:setup}

We adopt notation compatible with the two companion papers. A network induces a
representation map $x \mapsto \R$, with $\R \in \mathbb{R}^{d}$. An
\textbf{interpretability method} $\mathcal{M}$ is characterised by three
ingredients: (i) a \emph{hypothesis class} $\hypclass_{\complexity}$ indexed by a
scalar \emph{complexity knob} $\complexity \in \mathbb{R}_{+}$; (ii) a per-example
\emph{fit} functional $\fit(h; \cdot) \in [0,1]$ measuring how well a hypothesis
$h \in \hypclass_{\complexity}$ reproduces a target on one example; and (iii) a
\emph{selection rule} that, given a finite sample $\mathcal{D}_n$ of size $n$,
returns $\hat h = \arg\max_{h \in \hypclass_{\complexity}}
\widehat{\fit}(h; \mathcal{D}_n)$ and reports the achieved fit
$s(\mathcal{M}) = \widehat{\fit}(\hat h; \mathcal{D}_n)$.

The two methods instantiate this template as follows.

\paragraph{Diagnostic probing.} $\hypclass_{\complexity}$ is a family of probes
$q_{\bm\theta}(t \mid \mathbf r)$---e.g.\ multi-layer perceptrons whose depth or
width grows with $\complexity$, as in \citet{pimentel2020information}. The targets
are property labels $t$, the inputs are activations $\mathbf r$, and
$\widehat{\fit}$ is accuracy (or, in the information-theoretic view, the negative
cross-entropy that lower-bounds $\mathrm{I}(\T;\R)$). The dataset is
$\mathcal{D}_n = \{(\mathbf r_i, t_i)\}_{i=1}^n$.

\paragraph{Causal probing (DAS).} $\hypclass_{\complexity}$ is a family of bijective
alignment maps $\phi$---e.g.\ reversible networks
\citep{gomez2017reversible} whose depth $L_{\mathrm{rn}}$ or hidden size
$d_{\mathrm{rn}}$ grows with $\complexity$, as in \citet{sutter2025nonlinear}. The
``label'' of an example is the counterfactual output the \emph{target algorithm}
prescribes, the inputs are (base, source) pairs, and $\widehat{\fit}$ is the IIA:
the rate at which the \emph{frozen} network, intervened upon in the aligned basis,
reproduces that output. The dataset is
$\mathcal{D}_n = \{(\mathbf x^{\mathrm{base}}_i, \mathbf x^{\mathrm{src}}_i,
\mathbf y^{\mathrm{cf}}_i)\}_{i=1}^n$.

The crucial structural difference---which we will argue is the heart of the
matter---is that for DAS the hypothesis is not a classifier of $\mathbf r$. It is a
\emph{reparametrisation} $\phi$ of the hidden space, after which a \emph{fixed}
intervention and the \emph{fixed} downstream network produce the prediction. The
induced labelling function is $x \mapsto \mathrm{N}\!\big(\mathrm{intervene}_\phi(x)\big)$,
and the family is constrained to invertible $\phi$.

\section{A unified capacity measure}\label{sec:measure}

\subsection{Spuriousness as fit to randomised targets}

The propensity to ``find what you want'' is, statistically, the gap between the
fit a method reports and the fit it \emph{deserves}: the optimism of an
empirical-risk maximiser. Uniform-convergence theory names the leading term of this
optimism the \emph{Rademacher complexity} of the hypothesis class
\citep{bartlett2002rademacher}. We take this as our definition of spurious-fitting
capacity, expressed directly in the accuracy space the two methods report.

\begin{definition}[Spurious-fitting capacity]\label{def:capacity}
Fix a representation sample of size $n$ and let $\bm\sigma$ denote a target
\emph{randomisation}: a relabelling of the examples that destroys any genuine
relationship between inputs and targets while preserving their marginals. The
spurious-fitting capacity of a method $\mathcal{M}$ at complexity $\complexity$ is
\begin{equation}\label{eq:capacity}
  \rade_n(\hypclass_{\complexity})
    \;\coloneqq\;
    \Exp_{\bm\sigma}\!\left[\,
      \sup_{h \in \hypclass_{\complexity}}
      \widehat{\fit}\big(h;\, \bm\sigma(\mathcal{D}_n)\big)
    \,\right] \;-\; s_{\mathrm{chance}},
\end{equation}
i.e.\ the expected best fit the class achieves on randomised targets, centred at
chance performance.
\end{definition}

\autoref{eq:capacity} is, up to a scaling constant and the usual move from a
margin loss to $0/1$ accuracy, the empirical Rademacher complexity of
$\hypclass_{\complexity}$ on the representation sample. Its appeal is threefold.
First, it is \emph{method-agnostic}: it specialises to diagnostic probing (random
token labels) and to DAS (a random target algorithm, or equivalently permuted
counterfactual outputs) without changing form. Second, it is exactly the
quantity practitioners already estimate: it is the control-task accuracy of
\citet{hewitt2019control}, and the random-network alignment of
\citet{sutter2025nonlinear} is its causal analogue.\footnote{We deliberately phrase
\autoref{def:capacity} in accuracy space rather than as the textbook
$\Exp_{\bm\sigma}[\sup_h \frac1n \sum_i \sigma_i h(\mathbf r_i)]$ so that it is
directly comparable across the two methods and directly estimable from the
existing code paths, which already train on randomised labels.} Third, it is
\emph{cheap}: a Monte-Carlo estimate requires only re-running the existing fitting
pipeline on permuted targets.

\subsection{Two failure modes, kept separate}\label{sec:two-modes}

It is essential not to conflate \autoref{def:capacity} with the failure the two
theorems describe. They live at different scales.

\begin{description}[leftmargin=1.4em]
  \item[(a) Population vacuity.] As $\complexity \to \infty$, the
  \emph{population} estimand becomes uninformative: an arbitrarily powerful probe
  saturates the mutual-information bound \citep[Thm.~1]{pimentel2020information},
  and an arbitrarily complex alignment map achieves perfect IIA on any network
  \citep[Thm.~1]{sutter2025nonlinear}. In this limit \emph{both methods are equally
  doomed}; the limit itself does not distinguish them.
  \item[(b) Finite-sample optimism.] At finite $n$, the reported score exceeds the
  population value by an amount controlled by $\rade_n(\hypclass_{\complexity})$.
  This term \emph{does} distinguish the methods, and it is the only thing that can.
\end{description}

We therefore argue that the scientifically productive question is not the
$\complexity\to\infty$ limit---which the theorems have settled---but the
\emph{rate} at which $\rade_n(\hypclass_{\complexity})$ grows as we approach it.
``Failing gracefully'' is a statement about this rate. This separation also
disarms an obvious objection (``your own theorems say the limit is degenerate, so
why measure capacity?''): we are not studying the limit, we are studying the
finite-sample road to it.

\subsection{What ``graceful'' should mean, and the comparability problem}

A natural temptation is to compare the two methods by plotting their capacity
against a shared notion of ``size''---say, parameter count---and asking whether one
grows like $\bigO(\sqrt{P})$ where the other grows like $\bigO(P)$. We caution
against taking the exponent itself as the contribution, for two reasons. First,
generic norm- or parameter-based bounds yield the \emph{same} $\bigO(\sqrt{P/n})$
form for any reasonable class \citep{neyshabur2015norm}, so an a priori difference
in exponent is not expected; a measured difference would have to come from
structure, not from counting. Second, ``depth of an MLP probe'' and ``depth of a
RevNet alignment map'' are not the same unit, so a raw $\complexity$-axis is not
commensurable.

Both problems are solved by using \emph{capacity itself} as the common axis.
Concretely, we propose to plot, for each complexity level, the \textbf{signal}
(fit on \emph{correct} targets) against the \textbf{spurious capacity} (fit on
\emph{randomised} targets, \autoref{def:capacity}). Each method traces a curve in
this \emph{signal--spurious plane} as $\complexity$ varies. A method
\textbf{fails gracefully} if its curve hugs the top-left---high signal recovered
per unit of spurious capacity spent---and \textbf{fails sharply} if signal and
spurious capacity rise together. This operationalises ``one method fails more
gracefully than the other'' without ever committing to incommensurable units, and
it directly generalises the selectivity of \citet{hewitt2019control} (which is the
single-point gap $\text{signal} - \text{spurious}$) into a \emph{scaling curve}.

\section{A mechanism for divergence}\label{sec:mechanism}

If the two methods do diverge in the signal--spurious plane, the explanation must
lie in structural constraints that one method imposes and the other does not. We
identify three constraints that are present in DAS but absent in diagnostic
probing, each of which should \emph{throttle} spurious capacity below what the raw
parameter count of $\phi$ would suggest:

\begin{enumerate}[leftmargin=1.6em,itemsep=2pt]
  \item \textbf{Invertibility.} The alignment map $\phi$ is constrained to be a
  bijection. A bijection cannot collapse or discard directions of the hidden space
  to manufacture a fit; it can only rotate and warp them. This is a genuine
  capacity restriction with no analogue in an unconstrained probe.
  \item \textbf{The frozen network.} The prediction is produced by the
  \emph{fixed} downstream network applied to an intervened, $\phi$-aligned hidden
  state. The map cannot directly emit a label; it can only reroute information that
  the frozen weights will then process. Spurious fits must survive this fixed
  computation.
  \item \textbf{Multi-intervention consistency.} A single $\phi$ must
  simultaneously satisfy the counterfactual constraints of \emph{several}
  intervention types (in the hierarchical equality task, the two sub-equalities
  and their conjunction). These are many coupled constraints, not one label, and
  jointly satisfying them on randomised targets is harder than fitting a single
  randomised label.
\end{enumerate}

Our working hypothesis is that these constraints make
$\rade_n^{\mathrm{DAS}}(\complexity)$ grow \emph{more slowly} than
$\rade_n^{\mathrm{DP}}(\complexity)$---i.e.\ that DAS, despite the alarming
$\complexity\to\infty$ result of \citet{sutter2025nonlinear}, may be the more
graceful failure at the finite capacities used in practice. We stress that this is
a hypothesis to be tested, not a claim; the constraints above could equally be
swamped by the sheer flexibility a non-linear $\phi$ buys, in which case DAS would
fail \emph{more} sharply. Either outcome is informative.

\section{Experimental protocol}\label{sec:protocol}

We instantiate the study on the \textbf{hierarchical equality task}, the toy
setting of \citet{sutter2025nonlinear}, for which an implementation already exists
in this repository (a $3$-layer MLP trained to decide whether two pairs of
vectors are ``both equal or both unequal'', together with linear-probe and
linear-DAS pipelines and \emph{randomised-label} controls). The task is ideal
because the ground-truth intermediate variables ($\textsc{wx}$, $\textsc{yz}$) are
known, so ``signal'' is unambiguous.

\paragraph{Complexity axes.} We add a complexity knob to each method.
For \emph{diagnostic probing}, we replace the single logistic-regression probe with
a family of MLP probes of increasing depth and width
($\complexity \in \{$linear, $1$-hidden-layer, $2$-hidden-layer, $\dots\}$), as in
\citet{pimentel2020information}. For \emph{DAS}, we replace the orthogonal
(rotation) alignment map with a family of RevNet maps of increasing depth
$L_{\mathrm{rn}}$ and hidden size $d_{\mathrm{rn}}$, as in
\citet{sutter2025nonlinear}. The data budget $n$ is held \emph{fixed} across the
sweep so that we isolate the effect of $\complexity$ (in contrast to the existing
sweeps, which vary $n$).

\paragraph{Measurements.} At each complexity level and for several seeds we record
(i) the \emph{signal}: fit on correct targets (probe accuracy / IIA on the trained
network); and (ii) the \emph{spurious capacity}
$\widehat{\rade}_n(\hypclass_{\complexity})$ of \autoref{def:capacity}, estimated
by re-fitting on randomised targets (random token labels for probing; a randomly
relabelled / random target algorithm for DAS). We additionally repeat the DAS
measurement on a \emph{randomly initialised} network, which makes the spurious
axis vivid: any IIA there is, by construction, spurious.

\paragraph{Analysis.} We (1) plot signal and spurious capacity against
$\complexity$ for each method; (2) plot the two methods together in the
\emph{signal--spurious plane} of \autoref{sec:measure}; and (3) fit a simple
growth model $\widehat{\rade}_n(\complexity) \approx a\,\complexity^{\,b}$ to
read off, and compare, the empirical growth exponents $b$ with confidence
intervals. The deliverable is a statement of the form ``per unit of spurious
capacity, method X recovers more genuine signal than method Y on this task,'' with
the supporting curves.

\paragraph{What would falsify the framing.} If signal and spurious capacity rise in
lockstep for \emph{both} methods (curves collapsing onto the diagonal of the
signal--spurious plane), then capacity carries no method-distinguishing
information and the unification, while tidy, would be practically empty. We regard
making this falsifiable as a feature.

\section{Relation to prior work, and a caveat on consistency}\label{sec:related}

The control-task literature is the obvious antecedent. \citet{hewitt2019control}
introduce \emph{selectivity}---task accuracy minus control-task accuracy---as a
criterion for choosing probes, and \citet{voita2020information} propose minimum
description length as a complexity-aware alternative. Our spurious-fitting capacity
is, at a single complexity level, exactly a control-task score; the novelty here is
(i) reading it as a Rademacher complexity, (ii) lifting it from diagnostic to
\emph{causal} probing, and (iii) studying its \emph{scaling} rather than its value
at one point.

We must, however, confront a tension with \citet{pimentel2020information}, which
argues \emph{against} selectivity as a probe-\emph{selection} criterion: penalising
a complex probe for its control-task accuracy artificially discards capacity that
yields a tighter information estimate. We do not retract that argument; we
repurpose the quantity. Selectivity is a poor \emph{objective} to optimise---one
should still pick the best-performing probe for estimation---but it is an excellent
\emph{descriptive diagnostic} of how much a method \emph{could} fool itself, and
its growth rate is a property of the method, not a knob to tune. Measuring
fragility is not the same as selecting against it.

Finally, a caveat on what \autoref{def:capacity} does and does not capture.
Rademacher complexity is a statement about \emph{finite-sample} memorisation
(failure mode (b) of \autoref{sec:two-modes}). It is a tractable \emph{correlate}
of, but not identical to, the \emph{population-level} vacuity that the two theorems
prove (failure mode (a)). A method could have vanishing spurious capacity at finite
$\complexity$ and still be asymptotically vacuous. We therefore present capacity as
a measurable proxy for one---arguably the more practically dangerous---axis of
spuriousness, not as a complete account of it.

\section{Conclusion}\label{sec:conclusion}

We have argued that the diagnostic critique of \citet{pimentel2020information} and
the causal critique of \citet{sutter2025nonlinear} are two views of one
phenomenon---an empirical-risk maximiser flattering itself in proportion to its
capacity---and that this phenomenon has a single, measurable name: the Rademacher
complexity of the method's hypothesis class on the representations. Separating
population vacuity (where the methods are equally doomed) from finite-sample
optimism (where they need not be), we proposed to compare diagnostic probing and
DAS in a signal--spurious plane, and identified invertibility, the frozen network,
and multi-intervention consistency as concrete structural reasons DAS \emph{might}
fail more gracefully. The experiment is small, the infrastructure largely exists,
and the outcome---whichever way it falls---tells us something we currently only
guess at: not whether our interpretability tools can fool us, but how quickly.

\small
\bibliographystyle{plainnat}
\bibliography{references}

\end{document}
